\chapter{Local area and diameter estimates from Willmore energy}\label{sec:willmore-local-area-estimates}
Throughout this section $F:M^2\to\mathbb R^N$ is a smooth compact immersion of a nonempty surface without boundary. Set $X=F-y$, $r=|X|$, $M_s=\{r<s\}$, $A(s)=\mu(M_s)$ and $\mathcal W(E)=\frac14\int_E|\vec H|^2d\mu$. The cutoff calculation follows \cite{SimW93}, pp.~284--286. We keep the annular energy and the two boundary-radius terms explicit before deriving inequalities.

\begin{proposition}[The two-radius square identity]\label{prop:willmore-two-radius-identity}
For positive regular values $0<\sigma<\rho$ of $r$, let $M_{\sigma,\rho}=\{\sigma<r<\rho\}$. Then
\[
 \frac{A(\sigma)}{\sigma^2}
 +\int_{M_{\sigma,\rho}}\left|\frac{\vec H}4+\frac{X^\perp}{r^2}\right|^2d\mu
 =\frac{A(\rho)}{\rho^2}+\frac14\mathcal W(M_{\sigma,\rho})
 +\frac1{2\rho^2}\int_{M_\rho}X\cdot\vec H\,d\mu
 -\frac1{2\sigma^2}\int_{M_\sigma}X\cdot\vec H\,d\mu.
\]
There is no contribution from the intrinsic boundary of $M$, since it is empty.
\end{proposition}
\begin{proof}
Use the compactly supported Lipschitz ambient vector field
\[
 V(X)=\bigl(\max(r,\sigma)^{-2}-\rho^{-2}\bigr)_+X
\]
in the first variation identity $\int\operatorname{div}_MV=-\int\vec H\cdot V$. Its nonsmooth radii can be approximated by smooth radial cutoffs; the coefficients and tangential derivatives are uniformly bounded for fixed $\sigma>0$, and the two regular level sets have zero surface measure. Dominated convergence therefore justifies the Lipschitz test.

On $M_\sigma$ its tangential divergence is $2(\sigma^{-2}-\rho^{-2})$. On the annulus it is $2|X^\perp|^2/r^4-2\rho^{-2}$, and outside $M_\rho$ it is zero. Substitution and division by two give
\[
 \frac{A(\sigma)}{\sigma^2}+\int_{M_{\sigma,\rho}}
 \left(\frac{|X^\perp|^2}{r^4}+\frac{X\cdot\vec H}{2r^2}\right)d\mu
 =\frac{A(\rho)}{\rho^2}+\frac1{2\rho^2}\int_{M_\rho}X\cdot\vec H
 -\frac1{2\sigma^2}\int_{M_\sigma}X\cdot\vec H.
\]
Completing the square inside the annulus adds $|\vec H|^2/16$ on that annulus, exactly $\mathcal W(M_{\sigma,\rho})/4$ on the right. This explains why the energy in the exact two-radius formula is annular, before any enlargement used in an estimate.
\end{proof}

\begin{proposition}[A corrected area ratio is monotone]\label{prop:willmore-corrected-area-monotonicity}
For $s>0$ define
\[
 \Phi_y(s)=\frac{A(s)}{s^2}+\frac1{2s^2}\int_{M_s}X\cdot\vec H\,d\mu
 +\frac14\mathcal W(M_s).
\]
Then $\Phi_y$ is nondecreasing and
\[
 \Phi_y(s)=\pi m+\int_{M_s\setminus F^{-1}(y)}
 \left|\frac{\vec H}4+\frac{X^\perp}{r^2}\right|^2d\mu,
 \qquad m=\#F^{-1}(y).
\]
In particular $\Phi_y(0+)=\pi m$ and $\Phi_y(+\infty)=\mathcal W(F)/4$. The uncorrected area ratio need not be monotone for a nonminimal surface.
\end{proposition}
\begin{proof}
Move the inner curvature-radius term in Proposition \ref{prop:willmore-two-radius-identity} to the left and add $\mathcal W(M_\sigma)/4$ to both sides. The result is $\Phi_y(\rho)-\Phi_y(\sigma)$ equal to the nonnegative square integral over the annulus. The graph calculation of Lemma \ref{lem:willmore-puncture-flux} gives $A(\sigma)/\sigma^2\to\pi m$. It also gives $X^\perp=O(r^2)$ near each preimage, whence $\sigma^{-2}\int_{M_\sigma}X\cdot\vec H=O(\sigma^2)$; the inner energy tends to zero. If the pole misses the image, all three terms vanish for small $\sigma$. Letting $\sigma\downarrow0$ proves the formula for regular $s$.

For any other $s>0$, choose regular values increasing to $s$. The corresponding open sublevel sets increase to $M_s$. Monotone convergence for the nonnegative terms, and dominated convergence for the bounded signed curvature term, yield the same formula. It makes monotonicity immediate at every radius. Finally compactness of $M$ makes both radius-denominator terms tend to zero at infinity, while the energy tends to its full value. The correction terms involving mean curvature are essential in this computation.
\end{proof}

\begin{proposition}[A local comparison with an explicit constant]\label{prop:willmore-local-area-comparison}
For $0<\sigma<\rho$,
\[
 \frac{A(\sigma)}{\sigma^2}\leq3\frac{A(\rho)}{\rho^2}+\frac52\mathcal W(M_\rho).
\]
If $y$ has $m$ preimages, then
\[
 \pi m\leq3\frac{A(\rho)}{\rho^2}+\frac52\mathcal W(M_\rho).
\]
These estimates do not assume embeddedness or a pointwise bound on mean curvature.
\end{proposition}
\begin{proof}
First use regular radii in Proposition \ref{prop:willmore-two-radius-identity} and discard its square term. Put $a_s=A(s)/s^2$ and $E=\int_{M_\rho}|\vec H|^2$. The two radius terms are bounded in absolute value by $\tfrac12\sqrt{a_\rho E}$ and $\tfrac12\sqrt{a_\sigma E}$, using $r\leq s$ on $M_s$ and Cauchy--Schwarz. The annular energy term is at most $E/16$. Since
\[
 \tfrac12\sqrt{aE}\leq\tfrac12a+\tfrac18E,
\]
we obtain $a_\sigma\leq a_\rho+E/16+(a_\rho+a_\sigma)/2+E/4$. Rearrangement gives $a_\sigma\leq3a_\rho+5E/8$, which is the claim. Approximation from below by regular radii gives the formula at arbitrary radii; open sublevel sets and their integrals converge from below. Letting $\sigma\downarrow0$ gives the multiplicity consequence by the local sheet area expansion.
\end{proof}

\begin{theorem}[Global energy controls every local area ratio]\label{thm:willmore-uniform-quadratic-area}
For every $y\in\mathbb R^N$ and $s>0$,
\[
 \mu\bigl(F^{-1}(B_s(y))\bigr)\leq\mathcal W(F)s^2.
\]
The constant is independent of the total area, diameter and number of components. This is a consequence of the radial first-variation mechanism of \cite{SimW93}, pp.~284--285; the following proof retains separate inside and outside energies.
\end{theorem}
\begin{proof}
Fix a regular value $s$ and let $\rho\to\infty$ in Proposition \ref{prop:willmore-two-radius-identity}. Compactness eliminates the outer radius terms. Let $a=A(s)/s^2$, $w=\mathcal W(M_s)$ and $W=\mathcal W(F)$. Dropping the exterior square term gives
\[
 a\leq\frac{W-w}{4}-\frac1{2s^2}\int_{M_s}X\cdot\vec H
 \leq\frac{W-w}{4}+\sqrt{aw}.
\]
The last estimate again uses $r\leq s$ and Cauchy--Schwarz. Consequently
\[
 \left(\sqrt a-\frac{\sqrt w}{2}\right)^2\leq\frac W4,
 \qquad \sqrt a\leq\frac{\sqrt w+\sqrt W}{2}\leq\sqrt W.
\]
This proves $a\leq W$. For arbitrary $s$, take regular values increasing to $s$ and use convergence of the open-ball area. No limit surface or weak mean-curvature theory enters the argument.
\end{proof}

\begin{corollary}[Willmore energy gives a Gaussian-area bound]\label{cor:willmore-gaussian-area-bound}
Every smooth closed immersed surface in Euclidean space satisfies, for all centers $y$ and scales $t>0$,
\[
 F_{y,t}(M)=\frac1{4\pi t}\int_Me^{-|F-y|^2/(4t)}d\mu
 \leq\frac{\mathcal W(F)}\pi.
\]
Thus the supremum over centers and scales obeys the same bound. This estimate is not asserted to be sharp.
\end{corollary}
\begin{proof}
For $r\geq0$, write $e^{-r^2/(4t)}=\int_r^\infty\frac{s}{2t}e^{-s^2/(4t)}ds$. Tonelli's theorem gives
\[
 F_{y,t}=\frac1{8\pi t^2}\int_0^\infty s\,e^{-s^2/(4t)}A(s)\,ds.
\]
The choice of open or closed sublevel sets does not affect this integral in $s$. Apply Theorem \ref{thm:willmore-uniform-quadratic-area} and substitute $u=s^2/(4t)$. Since $\int_0^\infty s^3e^{-s^2/(4t)}ds=8t^2$, the result follows. Taking the supremum preserves the inequality. This connects an extrinsic bending-energy bound to the Gaussian functional used in mean curvature flow without claiming that Willmore energy is itself monotone along that flow.
\end{proof}

\begin{proposition}[A containing ball bounds area by energy]\label{prop:willmore-containing-ball}
If $R>0$ and the image lies in $\overline B_R(y)$, then $\mu(M)\leq R^2\mathcal W(F)$. Equality holds exactly when the immersion lies in $\partial B_R(y)$ and is minimal as a surface of that sphere.
\end{proposition}
\begin{proof}
Use the vector field $F-y$ in the closed-surface first variation. Its tangential divergence is two, so
\[
 2\mu(M)=-\int_M(F-y)\cdot\vec H
 \leq R\sqrt{\mu(M)}\left(\int_M|\vec H|^2\right)^{1/2}.
\]
This gives the inequality. Equality in the two estimates requires $|F-y|=R$ almost everywhere, hence everywhere, and $\vec H$ to be a fixed negative scalar multiple of $F-y$. The first-variation integral identifies that scalar as $2/R^2$. For a surface in a radius-$R$ sphere, the Euclidean and spherical mean curvature vectors satisfy $\vec H_{\mathbb R^N}=\vec H_{S^{N-1}(R)}-2(F-y)/R^2$. Therefore equality is equivalent to spherical minimality. Conversely that identity and $|F-y|=R$ give $\mathcal W=\mu(M)/R^2$ directly.
\end{proof}

\begin{theorem}[Area-energy bounds for extrinsic diameter]\label{thm:willmore-diameter-estimate}
If $M$ is connected and nonempty, its extrinsic diameter $D=\operatorname{diam}F(M)$ satisfies
\[
 \sqrt{\frac{\mu(M)}{\mathcal W(F)}}\leq D
 \leq\frac{29}{2\pi}\sqrt{\mu(M)\mathcal W(F)}.
\]
The upper constant is merely an explicit choice. The qualitative estimate is Lemma 1.1 of \cite{SimW93}, p.~283; the proof below also works with immersion multiplicity.
\end{theorem}
\begin{proof}
Fix a point of the image. The entire image lies in its closed radius-$D$ ball, so Proposition \ref{prop:willmore-containing-ball} gives the lower estimate. Write $A=\mu(M)$, $W=\mathcal W(F)>0$ and $\ell=\sqrt{A/W}$. Choose image points whose distance is $D$ and use one as center $y_0$. Connectedness makes the continuous distance function on the image attain every value in $[0,D]$. Hence for $j=0,\ldots,\lfloor D/\ell\rfloor$ choose image points $y_j$ with $|y_j-y_0|=j\ell$. Distinct points have pairwise distance at least $\ell$, so their open radius-$\ell/2$ balls are disjoint.

Each $y_j$ has at least one preimage. Apply Proposition \ref{prop:willmore-local-area-comparison} with outer radius $\ell/2$ at that center and sum. Disjointness of the preimages of these balls gives
\[
 \pi\bigl(1+\lfloor D/\ell\rfloor\bigr)
 \leq\frac{12A}{\ell^2}+\frac52W=\frac{29}{2}W.
\]
Since $D/\ell<1+\lfloor D/\ell\rfloor$, the upper bound follows, with a non-strict inequality as stated. This is an extrinsic diameter bound; it does not control intrinsic distances along a surface with long narrow geometry.
\end{proof}
